
\vspace{-0.8cm}

This chapter discusses how the \textbf{Automated Car Parking System} satisfies engineering standards, design constraints, Complex Engineering Problem (CEP) attributes, Complex Engineering Activity (CEA) attributes, and Knowledge Profile requirements. Since the project combines hardware, software, robotics, and IoT communication, safety, reliability, privacy, sustainability, and professional responsibility were considered during design and implementation.

\section{Compliance with Standards and Professional Practice}

The project followed basic engineering practices for embedded system design, hardware interfacing, modular programming, IoT communication, documentation, and testing. Relevant documentation such as Arduino UNO, Servo library, MFRC522 RFID module, ESP8266, and Firebase documentation was considered during development \cite{b8,b9,b10,b11,b12}.

\subsection{Software and Hardware Practice}

The software was developed using Arduino C/C++ in Arduino IDE. The code was organized into separate functions for RFID reading, gate control, LCD update, Wi-Fi connection, timestamp generation, and data transmission. Standard libraries such as \texttt{Servo.h}, \texttt{SPI.h}, \texttt{MFRC522.h}, \texttt{Wire.h}, and \texttt{LiquidCrystal\_I2C.h} were used.

The hardware design followed basic low-voltage embedded system interfacing practice. Arduino UNO was used as the controller, IR sensors were used for detection, MFRC522 used SPI communication, LCD used I2C communication, ESP-01S used serial communication, and the servo motor controlled the two-link robotic arm. Safe testing was performed using a small-scale prototype instead of a real vehicle gate.

\subsection{Professional Considerations}

The system stores RFID UID, entry/exit event, owner information, and timestamp. Therefore, privacy and responsible data use were considered. In real deployment, Wi-Fi credentials, database keys, and user records must be protected and should not be published in source code or reports.

\section{Design Constraints}

The project was affected by cost, component availability, sensor accuracy, power stability, mechanical alignment, and Wi-Fi reliability. Since it was developed as an academic prototype, low-cost components were selected while keeping the main functions operational.

\subsection{Ethical, Safety, and Sustainability Considerations}

The testing results and limitations were reported honestly. Related works and technical documentation were cited properly. RFID and online data should be handled responsibly to protect user privacy.

For safety, low-voltage components were used and the robotic arm was tested in a controlled lab environment. The servo movement was kept smooth to reduce sudden motion, and direct contact with the moving arm was avoided during operation.

From an environmental perspective, the system uses low-power and reusable electronic components. In practical use, automated parking can reduce unnecessary waiting time and vehicle movement near the gate, which may help reduce fuel consumption and emissions.

\section{Complex Engineering Problem (CEP) Coverage}

The project addresses a complex engineering problem because it integrates embedded control, sensors, authentication, robotic actuation, IoT communication, and real-time decision making.

\begin{table}[H]
\centering
\small
\caption{Mapping with Complex Engineering Problems (Ps)}
\label{tab:cep_mapping}
\resizebox{\textwidth}{!}{
\begin{tabular}{|p{0.07\textwidth}|p{0.24\textwidth}|p{0.5\textwidth}|p{0.17\textwidth}|}
\hline
\textbf{Ps} & \textbf{Attributes} & \textbf{How Ps are addressed} & \textbf{Evidence} \\
\hline
P1 & Depth of Knowledge & Applies embedded systems, sensors, RFID, servo control, IoT, and database communication. & Ch. 2, 3, 4 \\
\hline
P2 & Conflicting Requirements & Balances low cost, real-time response, reliable detection, smooth gate movement, and online data transfer. & Ch. 3, 5 \\
\hline
P3 & Depth of Analysis & Analyzes sensor response, RFID validation, slot counting, robotic arm motion, and Wi-Fi transmission. & Ch. 4, 5 \\
\hline
P4 & Familiarity of Issues & Parking automation is known, but robotic arm gate with IoT monitoring creates project-specific challenges. & Ch. 1, 2, 3 \\
\hline
P5 & Applicable Standards & Uses standard Arduino libraries, hardware interfacing practices, and IoT data handling principles. & Sec. 7.1 \\
\hline
P6 & Stakeholder Needs & Supports drivers, parking managers, and maintainers through automation and monitoring. & Ch. 3, Appendix \\
\hline
P7 & Interdependence & Sensors, RFID, servo, LCD, Wi-Fi, and server are interconnected; failure of one module affects the system. & Ch. 3, 4, 5 \\
\hline
\end{tabular}
}
\normalsize
\end{table}

\section{Complex Engineering Activities (CEA) Coverage}

The project involves complex engineering activities because it required multiple resources, hardware-software interaction, implementation, testing, and evaluation.

\begin{table}[H]
\centering
\small
\caption{Mapping of Complex Engineering Activities (As)}
\label{tab:activities_mapping}
\resizebox{\textwidth}{!}{
\begin{tabular}{|p{0.07\textwidth}|p{0.24\textwidth}|p{0.5\textwidth}|p{0.17\textwidth}|}
\hline
\textbf{As} & \textbf{Attributes} & \textbf{How As are addressed} & \textbf{Evidence} \\
\hline
A1 & Range of Resources & Uses Arduino UNO, IR sensors, RFID, ESP-01S, LCD, servo motor, robotic arm materials, libraries, and online database. & Ch. 3, 4 \\
\hline
A2 & Level of Interaction & Multiple modules interact, including sensors, RFID, servo, LCD, Wi-Fi module, and server. & Ch. 3, 4 \\
\hline
A3 & Innovation & Uses a two-link robotic arm gate instead of a simple static barrier. & Ch. 3, 5 \\
\hline
A4 & Society and Environment & Can reduce waiting time, improve parking management, reduce manual effort, and support smart city applications. & Ch. 1, 7 \\
\hline
A5 & Familiarity & Combines known technologies into a practical robotics-IoT parking automation prototype. & Ch. 2, 4 \\
\hline
\end{tabular}
}
\normalsize
\end{table}

\section{Knowledge Profile Mapping}

\begin{table}[H]
\centering
\small
\caption{Knowledge Profile Mapping}
\label{tab:knowledge_profile_mapping}
\resizebox{\textwidth}{!}{
\begin{tabular}{|p{0.09\textwidth}|p{0.35\textwidth}|p{0.48\textwidth}|}
\hline
\textbf{Profile} & \textbf{Knowledge Area} & \textbf{Project Evidence} \\
\hline
K1 & Natural science and engineering fundamentals & Sensor detection, electrical interfacing, and servo control. \\
\hline
K2 & Mathematics and computing & Control logic, slot counting, timestamp generation, and data formatting. \\
\hline
K3 & Engineering fundamentals & Embedded system design using Arduino and electronic modules. \\
\hline
K4 & Specialist knowledge & RFID authentication, IoT communication, and robotic arm gate control. \\
\hline
K5 & Engineering design & System architecture, circuit design, and hardware-software integration. \\
\hline
K6 & Engineering practice & Prototype development, testing, debugging, and documentation. \\
\hline
K7 & Social, environmental, and safety knowledge & Parking efficiency, safety, data privacy, and sustainability considerations. \\
\hline
K8 & Research literature and standards & Related works review and use of technical documentation. \\
\hline
\end{tabular}
}
\normalsize
\end{table}

\section{Summary}

This chapter presented the standards, design constraints, CEP coverage, CEA coverage, and Knowledge Profile mapping of the Automated Car Parking System. The project satisfies engineering requirements by integrating embedded control, robotics, IoT communication, and real-time parking management while considering safety, privacy, sustainability, and professional responsibility.

